% Cross-checks: scripts/comparison/.
\documentclass[11pt]{amsart}
\usepackage{amsmath,amssymb}
\setlength{\emergencystretch}{3em}
\tolerance=600
\usepackage{tikz-cd}
\usepackage{booktabs}
\usepackage[hidelinks]{hyperref}

\newtheorem{thm}{Theorem}[section]
\newtheorem{lem}[thm]{Lemma}
\newtheorem{cor}[thm]{Corollary}
\newtheorem{prop}[thm]{Proposition}
\theoremstyle{definition}
\newtheorem{defi}[thm]{Definition}
\newtheorem{rem}[thm]{Remark}
\newtheorem{ex}[thm]{Example}
\newtheorem{que}[thm]{Question}

\newcommand{\Same}{\mathrm{Same}}
\newcommand{\Top}{\mathrm{Top}}
\newcommand{\Mor}{\mathrm{Mor}}
\newcommand{\Iso}{\mathrm{Iso}}
\newcommand{\Presh}{\mathrm{Presh}}
\newcommand{\Sh}{\mathrm{Sh}}
\newcommand{\Sat}{\mathrm{Sat}}
\newcommand{\Loc}{\mathrm{Loc}}
\newcommand{\Hom}{\mathrm{Hom}}
\newcommand{\Map}{\mathrm{Map}}
\newcommand{\cosk}{\mathrm{cosk}}
\newcommand{\sk}{\mathrm{sk}}
\newcommand{\Fun}{\mathrm{Fun}}
\newcommand{\sS}{\mathcal S}
\newcommand{\E}{\mathcal E}
\newcommand{\tr}[1]{\tau_{\le #1}}
\newcommand{\re}{|{-}|}
\newcommand{\Wcyc}{W_{\mathrm{cyc}}}
\newcommand{\Wtest}{W_{\mathrm{test}}}
\newcommand{\Ob}{\mathrm{Ob}}
\newcommand{\Idem}{\mathrm{Idem}}
\newcommand{\sset}{\mathbf{sSet}}
\newcommand{\diag}{\mathrm{diag}}

\begin{document}

\title[Sheaves, truncations and reflections in the sameness lattice]{Sheaf equivalences, Postnikov truncations and preorder reflections in the lattice of two-out-of-three classes}
\author{Ryota Shibusawa}
\thanks{Corresponding author. E-mail: \texttt{r-shibusawa@daiichi-koudai.ac.jp}}
\address{Daiichi Institute of Technology}
\email{r-shibusawa@daiichi-koudai.ac.jp}
\keywords{two-out-of-three, localization, Grothendieck topology, subtopos,
  sheafification, Postnikov truncation, coskeleton, $\infty$-topos, lex modality,
  EI-category, preorder reflection, homomorphism order of graphs, finite monoid,
  tensor product of bisets}
\subjclass[2020]{18A32 (primary), 18F10, 18F20, 18N60, 55P60, 06B23, 20M12, 05C60 (secondary)}

\begin{abstract}
For a category $C$ let $\Same(C)$ be the complete lattice of classes of
morphisms that contain the identities and are closed under composition and
two-out-of-three. Every reflector $R$ contributes the element
$W_R=\{f:Rf\text{ invertible}\}$, and the question studied here is how the
lattice operations of $\Same(C)$ interact with the reflectors that arise in
sheaf theory, homotopy theory and order theory. Three results are proved.
(1) In the simplicial objects of any $\infty$-topos, the join of the
realization equivalences, the $(n{+}1)$-coskeletal equivalences and the
levelwise equivalences of a left-exact localization $L$ is exactly the class of
maps inverted by the $n$-truncated $L$-sheaf of the realization; for
simplicial sets this says that weak homotopy equivalences and coskeletal
equivalences generate the Postnikov $n$-equivalences under two-out-of-three.
(2) The sheaf equivalences $W_J$ of the Grothendieck topologies $J$ on a small
category $S$ form a meet-subsemilattice of $\Same(\Presh(S))$; joins are
preserved for finite acyclic categories and for finite categories whose two
topologies have an idempotent intersection of ideals, but not in general: on the
poset $\mathbb N$ the join of two sheaf classes is a proper class inverted by a
pro-limit. For finite categories the tower of bimodules representing the
alternating plus-construction tower is pro-isomorphic to the bimodule
representing the join sheafification, so that the plus tower has the join
sheafification as its colimit and this obstruction disappears; yet the tower
need not stabilize: an explicit monoid of order $12$
has a non-stabilizing tower, and a conjectured three-fold multiplicativity of
tensor products of ideals is refuted by the contracted monoid of a diamond
poset. (3) The morphisms $f\colon X\to Y$ admitting a morphism $Y\to X$ form the
largest element of $\Same(C)$ inverted by the preorder reflection; localizing
at it is the reflection into EI-categories, which for categories with binary
products or coproducts is the preorder reflection itself, so that the
EI-reflection of the category of graphs is the homomorphism preorder, whose
partial-order reflection is the homomorphism order.
\end{abstract}

\maketitle

\section{Introduction}

A \emph{sameness} on a category $C$ is a class $W$ of morphisms containing the
identities and closed under composition and two-out-of-three: if two of $f$,
$g$, $gf$ lie in $W$, so does the third. Weak equivalences of model categories,
the equivalences inverted by a reflector, the local isomorphisms of a
Grothendieck topology, and the morphisms with a morphism back are examples.
Samenesses form a complete lattice $\Same(C)$ under inclusion \cite{Same}, and
every functor $F\colon C\to D$ contributes the element
$F^*(\Iso)=\{f:Ff\text{ invertible}\}$, the largest sameness that $F$ inverts.
This paper studies the joins of such elements for the reflectors of sheaf
theory, homotopy theory and order theory. The join of two samenesses is the
smallest sameness containing both; it is in general much smaller than the class
of maps inverted by the composite reflector, and the results below identify
when the two coincide. Three theorems are proved.

\begin{thm}[Postnikov truncation as a join]\label{thm:I}
Let $\E$ be an $\infty$-topos, $\mathcal P=\Fun(\Delta^{\mathrm{op}},\E)$ its
simplicial objects, $\re$ the geometric realization, $\cosk_{n+1}$ the
coskeleton, $L$ a left-exact localization of $\E$ applied levelwise. In
$\Same(\mathcal P)$, for every $n\ge-1$,
\[
W_{\re}\vee W_{\cosk_{n+1}}\vee W_{L}=W_{L\tr n\re}.
\]
For simplicial sets: $\Wtest\vee W_{\cosk_{n+1}}=T_n$, where for $n\ge0$ $T_n$ is
the class of maps inducing isomorphisms on $\pi_i$ for $i\le n$ at all base
points, and $T_{-1}$ is the class of maps $X\to Y$ with $X=\varnothing$ iff
$Y=\varnothing$.
\end{thm}

\begin{thm}[topologies in the sameness lattice]\label{thm:II}
For a small category $S$ the assignment $J\mapsto W_J=\{f:a_Jf\text{ iso}\}$
is an order-embedding $\Top(S)\hookrightarrow\Same(\Presh(S))$ preserving all
meets. It preserves binary joins for finite acyclic categories and for finite
categories whose topologies have an idempotent intersection of ideals, but not
for the poset $\mathbb N$. For every finite category the tower of bimodules
representing the alternating plus-construction tower is pro-isomorphic to the
bimodule representing the join sheafification, so the plus tower has colimit
the join sheafification; the tower need not stabilize on the nose.
\end{thm}

\begin{thm}[the cycle class]\label{thm:III}
For every category $C$ the morphisms with a morphism back form the largest
sameness $\Wcyc$ inverted by the preorder reflection of $C$, and localizing at
$\Wcyc$ is the reflection into EI-categories. If $C$ has binary products or
binary coproducts, $C[\Wcyc^{-1}]$ is the reachability preorder of $C$; for
the category of graphs it is the homomorphism preorder, identifying finite
graphs with isomorphic cores.
\end{thm}

\subsection*{What is known.}
The lattice $\Same(C)$, its Galois functoriality and the comparison intervals
$[\Iso,F^*(\Iso)]$ are from \cite{Same}. The embedding of topologies, the
finite-sieve dictionary between topologies and idempotent ideals (for finite
$S$), and the simplicial-set case of Theorem~\ref{thm:I} are from
\cite{Postnikov}; meets, the two positive join theorems and the $\mathbb N$
counterexample from \cite{Topjoin}; the cycle class and its localization from
\cite{EI}; and the $\infty$-topos form of Theorem~\ref{thm:I} from
\cite{Infjoin}. The facts about localizations of $\infty$-topoi are in
\cite{HTT,ABFJ}; sheaf theory follows \cite{MLM,Elephant}; simplicial homotopy
theory \cite{GJ}; EI-categories \cite{Lueck}; the homomorphism order of graphs
\cite{HN}.

\subsection*{Related work.}
The lattice of Grothendieck topologies, or of subtoposes, is classical
\cite{MLM,Elephant}; its description by idempotent ideals for essential
localizations goes back to Kelly and Lawvere \cite{KL89}, and rigid topologies
on finite and Cauchy-complete categories, including the simplex category, are
characterized in \cite{Marques25}. The levels of the toposes of simplicial and
cubical sets, i.e.\ the coskeletal essential subtoposes, are determined in
\cite{KRRZ}. For $\infty$-topoi, Anel, Biedermann, Finster and Joyal describe
left-exact localizations by congruences and Grothendieck topologies
\cite{ABFJ,ABFJ2}. All of these concern the lattice of \emph{saturated}
classes, i.e.\ of localizations themselves. The focus here is different: the
bare two-out-of-three join of two such classes, before any saturation, and the
question of when it is already the class of a reflector. Theorems~\ref{thm:I}
and \ref{thm:II} answer it positively in the homotopical setting and for two
families of finite categories, and Section~\ref{sec:limits} shows what goes
wrong in general.

\subsection*{Provenance.}
Proposition~\ref{prop:lattice} is from \cite{Same}; Lemma~\ref{lem:loc},
Theorem~\ref{thm:join}, Theorem~\ref{thm:embed} (first part),
Lemma~\ref{lem:cosk} and Theorem~\ref{thm:sset} from \cite{Postnikov};
Theorem~\ref{thm:embed} (meets), Lemma~\ref{lem:reduction},
Theorems~\ref{thm:acyclic} and \ref{thm:monoid}, Lemma~\ref{lem:P} and
Theorem~\ref{thm:N} from \cite{Topjoin}; Proposition~\ref{prop:const} from
\cite{Infjoin}; and Section~\ref{sec:order} from \cite{EI}. Proofs of all these results are included, so that the paper is
self-contained; the remaining statements are new.

\subsection*{What is new.}
This paper supersedes and consolidates the preprints \cite{Same,Postnikov,Topjoin,EI,Infjoin}.
Besides the integration, the paper contains three results not in the cited
preprints: the three-fold multiplicativity conjecture of \cite[Conj.~6.12]{Topjoin}
is false (Proposition~\ref{prop:diamond}); for every finite category the
representing tower of the alternating plus construction is pro-isomorphic to
the representing object of the join sheafification, so that the plus tower has
the join sheafification as its colimit (Theorem~\ref{thm:essconst},
Corollary~\ref{cor:essconst}), which shows that the pro-limit
invariant of the $\mathbb N$ counterexample cannot occur for finite
categories; and there are finite monoids whose alternating tower does not
stabilize (Proposition~\ref{prop:cyclic}), so that
the finite join question is sharper than the stabilization criterion and remains
open. The proofs are self-contained.

\subsection*{Outline.}
Section~\ref{sec:lattice} sets up the lattice and the join theorem.
Sections~\ref{sec:sheaf}, \ref{sec:homotopy} and \ref{sec:order} prove
Theorems~\ref{thm:II}, \ref{thm:I} and \ref{thm:III}. Section~\ref{sec:limits}
collects the limits: the $\mathbb N$ example, the refuted conjecture, essential
constancy and the oscillating tower, and the open problems.

\section{The lattice of samenesses}\label{sec:lattice}

\begin{defi}
A \emph{sameness} of a category $C$ is a class $W\subseteq\Mor(C)$ containing
the identities, closed under composition, and satisfying two-out-of-three
\cite[Def.~2.1]{Same}. $\Same(C)$ is the set of samenesses ordered by inclusion.
For an $\infty$-category the same definition is applied to the morphisms of its
homotopy category. The isomorphisms form a sameness $\Iso(C)$, and every class
of the form $F^*(\Iso)$ below contains it; but a sameness need not contain
non-identity isomorphisms (Section~\ref{sec:order}).
\end{defi}

\begin{prop}\label{prop:lattice}
$\Same(C)$ is a complete lattice: meets are intersections, the top is $\Mor(C)$,
the bottom is $\{\mathrm{id}\}$, and the join of a family is the sameness generated by
its union. For a functor $F\colon C\to D$ and $W\in\Same(D)$ the class
$F^*(W)=\{f:Ff\in W\}$ is a sameness, and $F^*(\Iso)$ is the largest sameness
inverted by $F$; $F^*$ preserves meets and has a left adjoint $F_!$, so that
$\Same$ is a functor from $\mathbf{Cat}$ to complete lattices with Galois
connections.
\end{prop}
\begin{proof}
The three axioms are preserved by intersections, and $\Mor(C)$ satisfies them,
so $\Same(C)$ is a complete lattice with joins the meets of upper bounds. For
$F^*(W)$: $F$ preserves identities and composition, and if two of $f,g,gf$ lie
in $F^*(W)$ then two of $Ff,Fg,F(gf)$ lie in $W$, hence the third. A functor
inverts a sameness $V$ iff $V\subseteq F^*(\Iso)$. Meet preservation is
$F^*(\bigcap W_i)=\bigcap F^*(W_i)$, and a meet-preserving map of complete
lattices has a left adjoint.
\end{proof}

\begin{defi}\label{def:sat}
For $W\in\Same(C)$ an object $Z$ is \emph{$W$-local} if $\Hom(f,Z)$ is bijective
for all $f\in W$; $\Loc(W)$ is the class of local objects and
$\Sat(W):=\{f:\Hom(f,Z)\text{ bijective for all }Z\in\Loc(W)\}\supseteq W$ the
\emph{saturation}. $W$ is \emph{saturated} if $W=\Sat(W)$.
\end{defi}

\begin{lem}\label{lem:loc}
$\Loc(U\vee V)=\Loc(U)\cap\Loc(V)$. If $R$ is a reflector onto a full replete
subcategory $\mathcal R$ then $W_R:=R^*(\Iso)$ is saturated with
$\Loc(W_R)=\mathcal R$; consequently, if $\Loc(W)$ is reflective with reflector
$R$, then $\Sat(W)=W_R$.
\end{lem}
\begin{proof}
$\{f:\Hom(f,Z)\text{ bijective}\}$ is a sameness, so $Z$ is local for $U\vee V$
iff it is local for $U$ and for $V$. For a reflector with unit $\eta$: $Rf$ is
invertible iff $\Hom(Rf,Z)$ is bijective for all $Z\in\mathcal R$ iff
$\Hom(f,Z)$ is bijective for all $Z\in\mathcal R$, so
$W_R=\{f:\Hom(f,Z)\text{ bij.\ }\forall Z\in\mathcal R\}$ and
$\mathcal R\subseteq\Loc(W_R)$; if $Z$ is $W_R$-local then $\eta_Z\in W_R$ gives
a retraction $r$ of $\eta_Z$, and $\Hom(\eta_Z,RZ)$ bijective forces
$\eta_Zr=\mathrm{id}$, so $Z\in\mathcal R$. The last assertion is the
description of $W_R$ by orthogonality to $\mathcal R=\Loc(W)$.
\end{proof}

The saturation of Definition~\ref{def:sat} is orthogonality to the local
objects. It is in general larger than the class of maps inverted by the
localization functor $C\to C[W^{-1}]$: Theorem~\ref{thm:N} below gives a sameness
whose local objects are the terminal objects only, whose saturation in this
sense is therefore $\Mor$, but which is inverted by a limit functor that does not
invert every map.

\begin{thm}[join theorem]\label{thm:join}
Let $U,V,W\in\Same(C)$ with $W\subseteq U\vee V$, and let $P\colon C\to C$ be a
functor with a natural transformation $\pi\colon\mathrm{id}\Rightarrow P$ such
that (a) $\pi_X\in U\vee V$ for every $X$ and (b) $P(U)\cup P(V)\subseteq W$.
Then $U\vee V=P^*(W)$.
\end{thm}
\begin{proof}
$P^*(W)$ is a sameness containing $U$ and $V$ by (b), hence $U\vee V$. If
$Pf\in W$, the naturality square $\pi_Yf=Pf\,\pi_X$ has three of its four
sides in $U\vee V$, so $f\in U\vee V$ by two-out-of-three.
\end{proof}

\begin{cor}[two reflectors]\label{cor:two}
For reflectors $R_1,R_2$: $W_{R_1}\vee W_{R_2}=W_{R_2R_1}$ iff $R_2R_1$ inverts
$W_{R_2}$; in particular whenever $R_1$ preserves $R_2$-equivalences.
\end{cor}
\begin{proof}
Apply the theorem with $P=R_2R_1$, $W=\Iso$ and $\pi_X$ the composite of the
two units, a $W_{R_1}$-map followed by a $W_{R_2}$-map. Conversely the equality
gives $W_{R_2}\subseteq W_{R_2R_1}$.
\end{proof}

\section{Sheaf equivalences}\label{sec:sheaf}

Let $S$ be a small category, $\Top(S)$ the complete lattice of Grothendieck
topologies (ordered by inclusion of covering sieves), $a_J$ the sheafification
for $J$, a left-exact reflector onto $\Sh_J(S)$, and
$W_J:=a_J^*(\Iso)$, the class of $J$-local isomorphisms \cite[III.7]{MLM}.

\begin{thm}[{\cite[Thm.~3.1]{Postnikov}, \cite[Thm.~3.1]{Topjoin}}]\label{thm:embed}
$J\mapsto W_J$ is an order-embedding $\Top(S)\hookrightarrow\Same(\Presh(S))$
which preserves arbitrary meets: $W_{\bigwedge J_i}=\bigcap W_{J_i}$. Its image
consists of the saturated samenesses whose local objects form a subtopos. For
all $J,J'$, $W_{J\vee J'}=\Sat(W_J\vee W_{J'})$, so $W_J\vee W_{J'}\subseteq
W_{J\vee J'}$ with equality iff the join is saturated.
\end{thm}
\begin{proof}
$W_J$ is a sameness by Proposition~\ref{prop:lattice}, saturated with
$\Loc(W_J)=\Sh_J$ by Lemma~\ref{lem:loc}; the map is injective and monotone
because $J\le J'$ iff $\Sh_{J'}\subseteq\Sh_J$, and a saturated sameness is
determined by its local objects. For meets: a map $f$ is a $J$-local
isomorphism iff for every object $c$ and every $x\in X(c)$, $y\in Y(c)$ the
sieves $\{u: X(u)x=X(u)x'\}$ (for $x'$ with $fx=fx'$) and
$\{u: Y(u)y\in\mathrm{im}f\}$ are $J$-covering; these sieves do not depend on
$J$, and a sieve is covering for $\bigwedge J_i$ iff it is covering for every
$J_i$; hence $W_{\bigwedge J_i}=\bigcap W_{J_i}$. For the image: a saturated
sameness whose local objects are a subtopos $\Sh_J$ equals $W_J$ by
Lemma~\ref{lem:loc}. Finally $\Loc(W_J\vee W_{J'})=\Sh_J\cap\Sh_{J'}=\Sh_{J\vee J'}$,
so the saturation is $W_{J\vee J'}$.
\end{proof}

\subsection{The finite dictionary.}
Call a two-sided ideal $D$ of $S$ (a class of morphisms closed under
composition with arbitrary morphisms on both sides) \emph{idempotent} if every
member is a composite of two members. When $S$ has finitely many sieves on each
object, $\Top(S)\cong\Idem(S)^{\mathrm{op}}$ via $J\mapsto D_J=(\bigcap J(c))_c$
and $D\mapsto J_D(c)=\{R\supseteq D(c)\}$ \cite[Thm.~5.1]{Postnikov}. For a
finite monoid $M$ the topologies are thus the idempotent two-sided ideals under
reverse inclusion; for a finite acyclic category (no cycles, no non-identity
endomorphisms) they are the subsets $Y$ of objects, $D_Y$ the morphisms
factoring through $Y$, $\Sh_{J_Y}\simeq\Presh(Y)$ and
$W_Y=\{f: f_y\text{ bijective for }y\in Y\}$ \cite[\S5]{Topjoin}. The join
$J_D\vee J_E$ corresponds to the largest idempotent ideal $I'$ contained in
$D\cap E$; for a finite monoid $I'=(DE)^m$ for $m$ large.

\subsection{The alternating plus-construction tower.}
For topologies $J,J'$ write $\Phi_kX$ for the $k$-fold alternating plus
construction $X^{+J+J'+J\cdots}$, with units $u_k\colon X\to\Phi_kX$, each a
composite of maps of $W_J\cup W_{J'}$. For a finite monoid $M$ with idempotent
ideals $D,E$ (viewed as $M$--$M$-bisets), $X^{+D}=\Hom_M(D,X)$ and
$\Phi_kX=\Hom_M(T_k,X)$ with $T_k=D_k\otimes_M\cdots\otimes_MD_1$ the
alternating tensor tower ($D_1=D$, $D_2=E$, \dots), the units being induced by
the multiplications $\mu_k\colon T_k\to T_{k-1}$; and $a_DX=\Hom_M(D\otimes D,X)$.

\begin{lem}[reduction]\label{lem:reduction}
If for every $X$ there is $k=k(X)$ such that $\Phi_kX$ is a $(J\vee J')$-sheaf,
then $W_J\vee W_{J'}=W_{J\vee J'}$. For a finite monoid, the tower is
\emph{uniformly} stable, i.e.\ there is $k_0$ with $\Phi_kX\to\Phi_{k+1}X$
bijective for all $X$ and all $k\ge k_0$, iff $\mu_k$ is bijective for all
$k\ge k_0$; either condition implies join preservation.
\end{lem}
\begin{proof}
For $f\colon X\to X'$ in $W_{J\vee J'}$ take $k$ working for both; $\Phi_kf$ is a
map of $(J\vee J')$-sheaves inverted by $a_{J\vee J'}$, hence an isomorphism,
and $u_{X'}f=\Phi_kf\,u_X$ gives $f\in W_J\vee W_{J'}$. By Yoneda, $\Hom(\mu_{k+1},X)$
is bijective for all $X$ iff $\mu_{k+1}$ is; and if the units are eventually
bijective then $\Phi_{k_0}X$ is a $J$-sheaf and a $J'$-sheaf, hence a
$(J\vee J')$-sheaf, for every $X$.
\end{proof}

\begin{thm}[finite acyclic categories]\label{thm:acyclic}
For a finite acyclic category $S$ and all $Y,Y'\subseteq\Ob S$,
$W_Y\vee W_{Y'}=W_{Y\cap Y'}$; thus $\Top(S)\cong(2^{\Ob S})^{\mathrm{op}}$ is a
sublattice of $\Same(\Presh(S))$.
\end{thm}
\begin{proof}
Let $f\colon X\to X'$ be bijective on $Y\cap Y'$. Choose a linear extension
$p_1,\dots,p_m$ of the reachability preorder and put $S_i=\{p_1,\dots,p_i\}$.
Define the mixed presheaf $X_i$ by $X_i(q)=X'(q)$ for $q\in S_i$ and $X(q)$
otherwise, with restriction along $u\colon q\to q'$ equal to $X'(u)$, $X(u)$, or
$f_q\circ X(u)$ according as both, neither, or only $q$ lies in $S_i$ (the
remaining case cannot occur since $S_i$ is closed under predecessors);
functoriality is the naturality of $f$. The maps $\varphi_i\colon X_{i-1}\to X_i$
(identity away from $p_i$, $f_{p_i}$ there) are natural,
$f=\varphi_m\cdots\varphi_1$, and $\varphi_i$ is bijective away from $p_i$: it
lies in $W_Y$ if $p_i\notin Y$, in $W_{Y'}$ if $p_i\notin Y'$, and is an
isomorphism if $p_i\in Y\cap Y'$. Hence $f\in W_Y\vee W_{Y'}$; the reverse
inclusion is Theorem~\ref{thm:embed}.
\end{proof}

For finite monoids the argument is replaced by the tensor tower. Two facts about
a proper idempotent ideal $D$ of a finite monoid are used: every $d\in D$ can be
written $d=pgq$ with $p\in D$, $g\in D$ idempotent, $pg=p$ (write $d$ as a long
product of elements of $D$; two prefixes coincide, $p_i=p_iw$, and $g=w^\omega$),
and symmetrically $d=p'g'q'$ with $q'\in D$, $g'q'=q'$.

\begin{lem}[merging]\label{lem:P}
Let $X,Y,Z$ be two-sided ideals of a finite monoid with $Y$ idempotent, and
consider $X\otimes Y\otimes Z\otimes\cdots$ (further slots unaffected).
\begin{enumerate}
\item[(P)] If $c\in Y\cdot Z$, the class $[x,y,c,\dots]$ depends only on
$(xy,c,\dots)$. (P$'$) Mirror image: if $x\in X\cdot Y$, the class
$[x,y,c,\dots]$ depends only on $(x,yc,\dots)$.
\item[(A)] (Absorption) If the ideal of a slot is idempotent, an element $w$
of that slot may be replaced by $w'$ with $w=w'w''$, the factor $w''$ being
multiplied into the next slot to the right; so an interior slot may be assumed
to lie in the product of its ideal with the ideal to its right.
\item[(M)] If the conclusion of (P) holds for all tuples, then the
multiplication $X\otimes Y\otimes Z\otimes\cdots\to XY\otimes Z\otimes\cdots$ is
bijective; likewise for (P$'$) and $X\otimes YZ\otimes\cdots$.
\end{enumerate}
\end{lem}
\begin{proof}
(P) Write $c=y_1z_1$ with $y_1\in Y$, $z_1\in Z$, and $y_1=y'gy''$ with $g$
idempotent in $Y$, $y'g=y'$ (prefix factorization); moving $y'g$ to the left
and then $yy'$ to the left,
$[x,y,c]=[x,\,yy'g,\,y''z_1]=[xyy',\,g,\,y''z_1]$, each move legitimate because
the remaining entries lie in their slot ideals; the result depends on $xy$ and
on data determined by $c$. (P$'$) is the mirror image, with the suffix
factorization. (A) is the defining relation of the tensor product. (M) The
map is surjective; $[w,c,\dots]\mapsto[x,y,c,\dots]$ for any factorization
$w=xy$ is well defined by hypothesis and respects the relations of the target
(for $[wm,c]\sim[w,mc]$ use the factorization $wm=x(ym)$), so it is a section,
and $\mu$ is injective.
\end{proof}

\begin{thm}[idempotent intersections]\label{thm:monoid}
Let $M$ be a finite monoid and $D,E$ idempotent two-sided ideals with
$I:=D\cap E$ idempotent. Then $T_k\cong I\otimes_MI$ for $k\ge4$ by the
multiplication maps, $\Phi_4X\cong a_IX$ for every $M$-set $X$, and
$W_D\vee W_E=W_{D\cap E}$. The same holds for finite categories, with
profunctor tensors.
\end{thm}
\begin{proof}
Every product of two or more alternating factors $D,E$ lies between $I=I^2$
and $D\cap E=I$, hence equals $I$. Write the tuples of $T_k$ as
$[x_1,\dots,x_k]$ with slot ideals $D_k,\dots,D_1$ from left to right.

\emph{Step 1: merging slots $1,2$.} Let $t,t'$ be tuples differing only in
slots $1,2$ with $x_1x_2=x_1'x_2'$. Absorb slot $4$ into slot $3$ in both,
with the same factorization $x_4=pq$ (Lemma~\ref{lem:P}(A); here $k\ge4$ is
used): slot $3$ becomes $x_3p\in D_{k-2}D_{k-3}=I\subseteq D_{k-1}D_{k-2}$.
Lemma~\ref{lem:P}(P) for the slots $1,2,3$ shows that the classes of $t$ and
$t'$ agree, and (M) gives a bijection
$T_k\to I\otimes D_{k-2}\otimes\cdots\otimes D_1$.

\emph{Step 2: merging slots $2,3$ repeatedly.} In $I\otimes Y\otimes Z\otimes\cdots$
with $Y$ idempotent every element of the first slot lies in $I=I\cdot I\subseteq
I\cdot Y$, so (P$'$) applies to every tuple and (M) gives a bijection onto
$I\otimes YZ\otimes\cdots$; since $YZ=I$ and $I$ is idempotent this repeats
until two slots remain. All maps are multiplications, so the composite
$T_k\to I\otimes I$ is the canonical one and is an isomorphism of bisets.

Hence $\Phi_kX=\Hom_M(T_k,X)\cong\Hom_M(I\otimes I,X)=a_IX$, an $I$-sheaf, and
Lemma~\ref{lem:reduction} gives $W_D\vee W_E=W_I$.

\emph{Finite categories.} For a finite category $S$, ideals are
sub-profunctors $D\subseteq\Hom_S$, the tensor is the composite profunctor
(classes of composable chains modulo moving a morphism across a tensor sign),
$X^{+D}=\Hom(D,X)$ with $\Hom(P,X)(c)=\Hom_{\Presh(S)}(P(-,c),X)$, and
$a_DX=\Hom(D\otimes_SD,X)$. The factorization fact used in Lemma~\ref{lem:P}
holds for an idempotent ideal $Y$ of $S$: writing $y\in Y(a,c)$ as a composite
of more members of $Y$ than $S$ has morphisms, two partial composites with
domain $a$ coincide, $p_i=p_j=m\circ p_i$ with $m\in Y$ an endomorphism, and a
power $g=m^r$ is an idempotent of $Y$ with $g\circ p_i=p_i$; this gives
$y=y'\circ g\circ y''$ with $g\circ y''=y''$, and dually. The proofs of
Lemma~\ref{lem:P} and of Steps 1--2 use only the tensor relations, closure of
ideals under composition and this factorization, so they apply verbatim.
\end{proof}

The hypothesis fails already for the chain $a<b$ ($D_{\{a\}}\cap D_{\{b\}}$ is not
idempotent), which is covered by Theorem~\ref{thm:acyclic} instead. What
happens beyond both theorems is the subject of Section~\ref{sec:limits}.

\section{Postnikov truncations}\label{sec:homotopy}

\subsection{Coskeletal equivalences.}
For the simplex category $\Delta$ (and the cube category with faces and
degeneracies) the two-sided ideals are $\varnothing$, the ideals $D_k$ of
morphisms factoring through an object of dimension $\le k$, and $\Mor$; all are
idempotent, so the topologies form a chain and the $J_{D_k}$-sheaves are the
$k$-coskeletal presheaves \cite[Thm.~6.1]{Postnikov}. The sheaf samenesses are
the coskeletal equivalences
$W_{D_k}=\{f:\cosk_kf\text{ iso}\}=\{f: f_m\text{ bijective for }m\le k\}$,
a descending chain with $\bigcap_kW_{D_k}=\Iso$.

\begin{lem}\label{lem:cosk}
For a Kan complex $K$ and $n\ge-1$, $\cosk_{n+1}K$ is a Kan complex, the unit
$K\to\cosk_{n+1}K$ induces isomorphisms on $\pi_i$ for $i\le n$ at all base
points, and $\pi_i(\cosk_{n+1}K)=0$ for $i\ge n+1$ (for $n=-1$: empty or
contractible).
\end{lem}
\begin{proof}
An $m$-simplex of $\cosk_{n+1}K$ is a map $\sk_{n+1}\Delta^m\to K$. A horn
$\Lambda^k_m\to\cosk_{n+1}K$ is a map $\sk_{n+1}\Lambda^k_m\to K$; for $m\ge n+3$
this skeleton is $\sk_{n+1}\Delta^m$ and the horn is already a simplex, for
$m\le n+2$ it is a horn in $K$ whose filler restricts to the required simplex.
Homotopy groups of a Kan complex are computed from simplices of dimension
$i,i+1$ with boundary at the base point, where the unit is bijective for
$i\le n$; for $i\ge n+1$ every map $\partial\Delta^{i+1}\to\cosk_{n+1}K$ is a map
$\sk_{n+1}\Delta^{i+1}\to K$, i.e.\ extends to $\Delta^{i+1}$, so $\pi_i=0$.
\end{proof}

\begin{thm}\label{thm:sset}
In $\Same(\sset)$, with $\Wtest$ the weak equivalences,
$T_n=\{f:\pi_if\text{ iso for }i\le n\text{ at all base points}\}$ for $n\ge0$ and
$T_{-1}=\{f\colon X\to Y: X=\varnothing\Leftrightarrow Y=\varnothing\}$:
$\Wtest\vee W_{D_{n+1}}=T_n$ for all $n\ge-1$, and $\Wtest=\bigcap_n(\Wtest\vee W_{D_n})$.
\end{thm}
\begin{proof}
Apply Theorem~\ref{thm:join} with $P=\cosk_{n+1}\mathrm{Ex}^\infty$ and
$W=\Wtest$: $\pi_X$ is the Kan replacement followed by a coskeleton unit,
$P$ preserves weak equivalences (Lemma~\ref{lem:cosk} and Whitehead's theorem
for Kan complexes), and for $f\in W_{D_{n+1}}$, $\mathrm{Ex}^\infty f$ is again
bijective in dimensions $\le n+1$ (as $\mathrm{Ex}^\infty$ is computed
dimensionwise from the same low-dimensional data), so $Pf$ is a weak equivalence
by Lemma~\ref{lem:cosk}. Thus the join is $P^*(\Wtest)=T_n$ by
Lemma~\ref{lem:cosk}. The intersection statement is Whitehead's theorem.
\end{proof}

\subsection{Simplicial objects of an $\infty$-topos.}
Let $\E$ be an $\infty$-topos, $\mathcal P=\Fun(\Delta^{\mathrm{op}},\E)$,
$\re=\operatorname{colim}$, $c$ the constant functor (fully faithful, since
$\Delta^{\mathrm{op}}$ has an initial object), $\cosk_{n+1}$ the coskeleton
(a right Kan extension computed by finite limits, hence commuting with every
functor preserving finite limits), and $\tr n$ the truncation. $W_{\re}$ and
$W_{\cosk_{n+1}}$ are the saturated samenesses of the reflectors $c\re$ and
$\cosk_{n+1}$. A constant object $cK$ is $(n{+}1)$-coskeletal iff $K$ is
$n$-truncated, since the unit at level $n+2$ is $K\to K^{S^{n+1}}$
\cite[Lemma~8.2]{Topjoin}; but the levels $K^{|\sk_{n+1}\Delta^m|}$ of
$\cosk_{n+1}cK$ are not truncated. The key fact is that its realization is.

\begin{prop}\label{prop:const}
For $K\in\E$ and $n\ge-1$ the realization of the coskeleton unit
$cK\to\cosk_{n+1}cK$ is the truncation $K\to\tr nK$; so
$|\cosk_{n+1}cK|\simeq\tr nK$ naturally in $K$.
\end{prop}
\begin{proof}
For $\E=\sS$ choose a Kan complex $K'$ with $|K'|\simeq K$ and consider the
bisimplicial sets $S_{m,k}=\Hom(\Delta^m\times\Delta^k,K')$,
$T_{m,k}=\Hom(\sk_{n+1}\Delta^m\times\Delta^k,K')$ and
$T'_{m,k}=\Hom(\Delta^m\times\Delta^k,\cosk_{n+1}K')$ with the restriction maps
$S\to T\to T'$. For fixed $m$ the columns $S_{m,\bullet}=K'^{\Delta^m}$ and
$T_{m,\bullet}=K'^{\sk_{n+1}\Delta^m}$ are simplicial mapping spaces into a Kan
complex, of homotopy types $K$ and $\Map(|\sk_{n+1}\Delta^m|,K)$, so $S\to T$
models the coskeleton unit $cK\to\cosk_{n+1}cK$, and the realization of a
simplicial space obtained by realizing a bisimplicial set columnwise is the
diagonal of the bisimplicial set \cite[IV.1.9 and Ex.~IV.1.7]{GJ}; thus the
realized unit is $\diag S\to\diag T$. For fixed
$k$ the rows are $K'^{\Delta^k}\to\cosk_{n+1}(K'^{\Delta^k})\to(\cosk_{n+1}K')^{\Delta^k}$;
by Lemma~\ref{lem:cosk} the first map is an $n$-equivalence to an
$n$-truncated Kan complex, and so is the composite (the exponential of the unit
of $K'$), whence the second map is a weak equivalence. A map of bisimplicial sets that is a weak equivalence in each row induces a
weak equivalence of diagonals \cite[IV.1.7]{GJ}; applied to $T\to T'$ and to
the rowwise equivalences $K'^{\Delta^k}\to K'$ and
$(\cosk_{n+1}K')^{\Delta^k}\to\cosk_{n+1}K'$ (evaluation at the vertex $0$ of
$\Delta^k$, which has a section and is a weak equivalence of Kan complexes),
this identifies $\diag S\simeq K'$ and $\diag T\simeq\diag T'\simeq\cosk_{n+1}K'$
compatibly with the units, so that $\diag S\to\diag T$ is the coskeleton unit
of $K'$, i.e.\ the Postnikov section by Lemma~\ref{lem:cosk}.
For a presheaf $\infty$-topos everything is computed objectwise, and for
$\E=L\,\Presh(C)$ the coskeleton is computed in $\Presh(C)$ while
$|{-}|_\E=L|{-}|$ and $L\tr n=\tr nL$ \cite[5.5.6.28]{HTT}, so the identification
is transported by $L$.
\end{proof}

\begin{proof}[Proof of Theorem~\ref{thm:I}]
First $W_{\re}\vee W_{\cosk_{n+1}}=W_{\tr n\re}$: apply Theorem~\ref{thm:join}
with $P=c\tr n\re$ and $\pi_X\colon X\to c|X|\to c\tr n|X|$. Condition (b):
$P$ inverts $W_{\re}$, and inverts $W_{\cosk_{n+1}}$ because the local objects
of the join are the constant objects on $n$-truncated objects, a reflective
subcategory with reflector $P$, so $P$ inverts the whole join. Condition (a):
the first map is the unit of $\re\dashv c$, in $W_{\re}$; the second is
$c(K\to\tr nK)$ for $K=|X|$, and by Proposition~\ref{prop:const} it is, up to
the equivalence $c|\cosk_{n+1}cK|\simeq c\tr nK$, the composite of the
coskeleton unit $cK\to\cosk_{n+1}cK$ (in $W_{\cosk_{n+1}}$) with the unit of
$\re\dashv c$ at $\cosk_{n+1}cK$ (in $W_{\re}$). Then apply
Corollary~\ref{cor:two} to the reflectors $c\tr n\re$ and levelwise $L$: their
composite inverts $W_L$ because $L$ commutes with truncation and with colimits.
For simplicial sets, Theorem~\ref{thm:sset}.
\end{proof}

\begin{figure}[ht]
\[
\begin{tikzcd}[row sep=normal, column sep=small]
& & W_{L\tr n\re} & & \\
& W_{\tr n\re}\arrow[ur,dash] & W_{L\re}\arrow[u,dash] & W_{\cosk_{n+1}L}\arrow[ul,dash] & \\
& W_{\re}\arrow[u,dash]\arrow[ur,dash] & W_{\cosk_{n+1}}\arrow[ul,dash]\arrow[ur,dash] & W_{L}\arrow[ul,dash]\arrow[u,dash] & \\
& & \Iso\arrow[ul,dash]\arrow[u,dash]\arrow[ur,dash] & &
\end{tikzcd}
\]
\caption{Theorem~\ref{thm:I}: the sublattice of $\Same(\Fun(\Delta^{\mathrm{op}},\E))$
generated by the three reflector classes; every join is again a reflector
class (lines are joins, read upwards).}\label{fig:cube}
\end{figure}

\section{The cycle class and preorder reflection}\label{sec:order}

Write $X\preceq Y$ if $\Hom(X,Y)\neq\varnothing$ (the reachability preorder of
$C$) and let $\Wcyc(C)=\{f\colon X\to Y:\ Y\preceq X\}$, the morphisms with a
morphism back. A category is an \emph{EI-category} if all its endomorphisms
are invertible \cite{Lueck}.

\begin{thm}\label{thm:cyc}
$\Wcyc(C)$ is a sameness, equal to $P^*(\Iso)$ for the preorder reflection
$P\colon C\to(\Ob C,\preceq)$; it is $\{\mathrm{id}\}$ iff $C$ is acyclic, $\Iso$ iff
$C$ is an EI-category, and $\Mor$ iff $\preceq$ is symmetric. The localization
$C[\Wcyc^{-1}]$ is an EI-category with the same reachability preorder, and
$C\to C[\Wcyc^{-1}]$ is universal among functors to EI-categories.
\end{thm}
\begin{proof}
$Pf$ is an isomorphism in the preorder iff $X\preceq Y\preceq X$ iff $f$ has a
back, so $\Wcyc=P^*(\Iso)$ is a sameness. The three extremes are immediate. For
the localization: a zigzag $X\to\cdots\to Y$ in $C[\Wcyc^{-1}]$ exists iff
$X\preceq Y$ (backward arrows are in $\Wcyc$, hence reversible in $\preceq$), so
the preorder is preserved; an endomorphism of $X$ in the localization is a
zigzag whose forward arrows all lie in strongly connected classes reachable from
and reaching $X$, so each forward arrow has a back and is inverted; hence the
zigzag is invertible and the localization is EI. Conversely a functor to an
EI-category $D$ sends $f\colon X\to Y$ with back $g$ to $F(f)$ with
$F(g)F(f)$ and $F(f)F(g)$ invertible endomorphisms, so $F(f)$ is invertible; so
$F$ inverts $\Wcyc$ and factors uniquely through the localization.
\end{proof}

\begin{thm}\label{thm:preorder}
If $C$ has binary coproducts or binary products, $C[\Wcyc^{-1}]$ is the
reachability preorder $(\Ob C,\preceq)$: the EI-reflection is the preorder
reflection.
\end{thm}
\begin{proof}
For parallel $g,g'\colon B\to C'$ let $j_1,j_2\colon B\to B\sqcup B$ be the
injections and $\nabla$ the fold map, with $\nabla j_1=\nabla j_2=\mathrm{id}$. Then
$\nabla$ has the back $j_1$, so it is inverted in $C[\Wcyc^{-1}]$, where
$j_1=\nabla^{-1}=j_2$ and $g=[g,g']j_1=[g,g']j_2=g'$. Parallel zigzags are then
equal too, so the localization is a preorder, and by Theorem~\ref{thm:cyc} it
is $\preceq$. With products use the projections and the diagonal.
\end{proof}

\begin{cor}\label{cor:graphs}
In a topos, and in every presheaf category, $\E[\Wcyc^{-1}]$ is the hom-preorder.
For directed graphs it is the homomorphism preorder, whose partial-order
reflection is the homomorphism order of \cite{HN}; two finite graphs are
identified iff their cores are isomorphic. An EI-category with binary
(co)products is a preorder; a groupoid with binary coproducts is equivalent to a
discrete category.
\end{cor}

\begin{prop}[natural samenesses]\label{prop:natural}
Call $C\mapsto W(C)$ natural if $F(W(C))\subseteq W(D)$ for all functors $F$.
Then $\{\mathrm{id}\}$, $\Iso$, $W_{\mathrm{obj}}$ (source isomorphic to target), $\Wcyc$,
$\Mor$, the endomorphism class $W_{\mathrm{end}}$ and its intersection with
$\Iso$ are natural; $\Iso$ and $W_{\mathrm{end}}$ are incomparable with join
$W_{\mathrm{obj}}$; and every natural sameness other than $\Mor$ is contained in
$\Wcyc$.
\end{prop}
\begin{proof}
Each listed class is preserved by every functor (functors preserve identities,
isomorphisms, endomorphisms and the existence of a morphism back), and each is
a sameness: for $W_{\mathrm{obj}}$ and $\Wcyc$ this is the object-determined
criterion (a class of the form $\{f\colon X\to Y: X\,R\,Y\}$ for an equivalence
relation $R$ on objects is a sameness), for $W_{\mathrm{end}}$ two-out-of-three
holds because a composite of two morphisms is an endomorphism iff both are,
when one of them is. An isomorphism between distinct objects lies in $\Iso$
but not in $W_{\mathrm{end}}$, and a non-invertible endomorphism in
$W_{\mathrm{end}}$ but not in $\Iso$, so the two are incomparable; both lie in
$W_{\mathrm{obj}}$, and conversely $f\colon X\to Y$ with an isomorphism
$\varphi\colon X\to Y$ is $\varphi\circ(\varphi^{-1}f)$ with $\varphi^{-1}f$ an
endomorphism, so $W_{\mathrm{obj}}=\Iso\vee W_{\mathrm{end}}$. Finally let $W$
be natural and $f\colon X\to Y$ in $W(C)$ with $Y\not\preceq X$; for any
morphism $u\colon A\to B$ of any category $D$ there is a functor $F$ from the
category generated by $f$ (an acyclic category with two objects, since
$Y\not\preceq X$ and the only endomorphisms are identities) to $D$ with
$F(f)=u$, so $u\in W(D)$ and $W=\Mor$. Hence a proper natural sameness
contains only morphisms with a back, i.e.\ $W\subseteq\Wcyc$.
\end{proof}

\begin{prop}\label{prop:toposcyc}
In a topos, $\Loc(\Wcyc)$ is the class of subterminal objects and
$\Sat(\Wcyc)$ the class of maps between objects with the same support; $\Wcyc$ is
saturated in $\mathbf{Set}$ and in $\sset$ but not in graphs.
\end{prop}
\begin{proof}
If $U$ is subterminal then $\Hom(X,U)$ has at most one element and is non-empty
iff $X\to U$ exists; for $f\colon X\to Y$ with a back, $X\to U$ exists iff
$Y\to U$ exists, so $\Hom(f,U)$ is bijective and $U$ is local. Conversely let
$Z$ be local and apply $\Hom(-,Z)$ to the projection $p_1\colon Z\times Z\to Z$,
which has the back $\Delta$: bijectivity means every map $Z\times Z\to Z$
factors through $p_1$; for $p_2=hp_1$ composing with $\Delta$ gives $h=\mathrm{id}$,
so $p_1=p_2$ and $Z\to1$ is a monomorphism. Thus $\Loc(\Wcyc)=\mathrm{Sub}(1)$.
A map $f\colon X\to Y$ is orthogonal to all subterminals iff for every
subterminal $U$ a map $X\to U$ exists iff a map $Y\to U$ exists, i.e.\ iff
$X$ and $Y$ have the same support (the support of $X$ being the least $U$ with
$X\to U$). Maps with a back preserve supports, so $\Wcyc\subseteq\Sat(\Wcyc)$,
with equality iff any two objects with the same support admit morphisms in
both directions: true in $\mathbf{Set}$ and in $\sset$ (all non-empty objects
have support $1$ and admit maps both ways), false for graphs (the graph with
one vertex and no edge and the graph with one vertex and a loop have the same
support, but there is no morphism from the looped graph to the loopless one).
\end{proof}

\section{Limits, counterexamples and open problems}\label{sec:limits}

\subsection{Joins are not preserved in general.}

\begin{thm}\label{thm:N}
On the poset $\mathbb N$ let $J_{\mathrm{ev}},J_{\mathrm{od}}$ be the topologies
of the ideals of morphisms passing through an even, resp.\ odd, number. Their
join is the degenerate topology, so $W_{J_{\mathrm{ev}}\vee J_{\mathrm{od}}}=\Mor$,
but $W_{\mathrm{ev}}\vee W_{\mathrm{od}}$ is contained in the class of maps
inverted by $\lim_{\mathbb N^{\mathrm{op}}}$, which does not contain the map from
the presheaf of final segments to the terminal presheaf.
\end{thm}
\begin{proof}
Degeneracy: $\varnothing$ covers $0$ for $J_{\mathrm{od}}$, and inductively the
least covering sieve of $m$ for the topology of the other parity consists of
arrows from smaller objects, which are covered by $\varnothing$; transitivity
gives $\varnothing\in J(m)$. $W_{\mathrm{ev}}=\{f:f_y\text{ bijective for even }y\}$
and the evens are cofinal, so $\lim_{\mathbb N^{\mathrm{op}}}$ inverts
$W_{\mathrm{ev}}$; likewise $W_{\mathrm{od}}$; hence it inverts their join. For
$X(n)=\{x\ge n\}$ the limit is empty, while the limit of the terminal presheaf
is a point.
\end{proof}

\subsection{A refuted conjecture.}
\cite[Conj.~6.12]{Topjoin} proposed that for idempotent ideals $X,Y,Z\ne M$ of a
finite monoid the multiplication $X\otimes Y\otimes Z\to XY\otimes Z$ is always
bijective, which by Lemma~\ref{lem:P} would give join preservation for all
finite monoids. It holds for all monoids of order $\le5$ but fails in general.

\begin{prop}\label{prop:diamond}
Let $P$ be the diamond poset $0<1,2<3$ (with $1,2$ incomparable) and $M$ its
contracted monoid: an identity, the nine morphisms $a\to b$ ($a\le b$), and a
zero, with product the composite when defined and $0$ otherwise. Let
$X=Y$ be the ideal of all morphisms except the identities of $0$ and $3$, and
$Z=M\,1_3\,M$, all proper idempotent ideals. Then
$[0\to1,\,1\to3]\ne[0\to2,\,2\to3]$ in $X\otimes_MY$ although both have product
$0\to3$, and the classes $[0\to1,1\to3,1_3]$ and $[0\to2,2\to3,1_3]$ remain
distinct in $X\otimes Y\otimes Z$, so $X\otimes Y\otimes Z\to XY\otimes Z$ is not
injective.
\end{prop}
\begin{proof}
Write $x=0\to1$, $y=1\to3$. A tensor move replaces $(x,y)$ by $(x',my)$ with
$x=x'm$, $x'\in X$, or by $(xm,y')$ with $y=my'$, $y'\in Y$. The factorizations
of $x$ in $M$ are $x\cdot1_1$ and $1_0\cdot x$ with $1_0\notin X$, and those of
$y$ are $1_1\cdot y$ and $y\cdot1_3$ with $1_3\notin Y$; so all moves are trivial
and the class of $(x,y)$ is the singleton $\{(x,y)\}$; likewise for
$(0\to2,2\to3)$. Let $Z=M\,1_3\,M=\{0\to3,1\to3,2\to3,1_3,0\}$, an idempotent
proper ideal, and $c=1_3$. In $X\otimes Y\otimes Z$ the additional moves act on
the last two slots: $(y,c)\mapsto(y',mc)$ with $y=y'm$, $y'\in Y$, and
$(y,c)\mapsto(ym,c')$ with $c=mc'$, $c'\in Z$. Since $1_1\in Y$, the
factorization $y=1_1\cdot y$ gives the move $(x,y,1_3)\mapsto(x,1_1,y)$, and
from $(x,1_1,y)$ the only non-trivial move is the inverse one; so the class of
$(x,y,c)$ is $\{(x,y,1_3),(x,1_1,y)\}$, and likewise the class of
$(0\to2,2\to3,1_3)$ is $\{(0\to2,2\to3,1_3),(0\to2,1_2,2\to3)\}$. The two classes
are disjoint but both map to $[0\to3,1_3]$ in $XY\otimes Z$.
\end{proof}

\begin{rem}
The conjecture was a sufficient condition applied to all triples of ideals,
whereas only the triples along the alternating pattern $D,E,D,E,\dots$ enter
the tower; a finite enumeration suggests that on this monoid the towers of all
pairs stabilize, so the conjecture is stronger than what join preservation
requires.
\end{rem}

\subsection{Essential constancy for finite categories.}
The obstruction in Theorem~\ref{thm:N} is a pro-limit invariant of the
alternating plus-construction tower. For finite categories no such invariant
exists. We use the Yoneda lemma for pro-objects: for a small category
$\mathcal A$, the functor $\mathrm{Pro}(\mathcal A)^{\mathrm{op}}\to\Fun(\mathcal A,\mathbf{Set})$,
$(T_k)_k\mapsto\operatorname{colim}_k\Hom(T_k,-)$, is fully faithful, so that
$\Hom_{\mathrm{Pro}(\mathcal A)}((T_k)_k,X)=\operatorname{colim}_k\Hom(T_k,X)$ for
$X\in\mathcal A$ and two pro-objects with naturally isomorphic functors are
isomorphic \cite[\S6.1]{KS}.

\begin{thm}\label{thm:essconst}
Let $M$ be a finite monoid, $D,E$ idempotent ideals, $I'$ the largest
idempotent ideal in $D\cap E$, and $m\ge1$ such that $(DE)^m=(ED)^m=I'$ (the
descending chains $(DE)^j$, $(ED)^j$ stabilize at idempotent ideals contained in
$D\cap E$, which contain $I'=I'^{2j}$, hence equal $I'$). Put
$P=I'\otimes_MI'$, so that $a_{I'}X=\Hom(P,X)$. For every $M$-set $X$,
$\operatorname{colim}_k\Phi_kX\cong a_{I'}X$ naturally; hence the pro-object
$(T_k,\mu_k)$ is pro-isomorphic to the constant object $P$, and for $k$ large
$a_{I'}X$ is a natural retract of $\Phi_kX$ compatibly with the units. Moreover
$\mathrm{im}(T_{k+l}\to T_k)=I'\cdot T_k$ for $l\ge2m$, so the tower stabilizes
($\mu_k$ bijective for large $k$) iff $T_k=I'\cdot T_k$ for large $k$, iff every
class of $T_k$ has a representative whose first slot lies in $I'$; in that case
$W_D\vee W_E=W_{I'}$.
\end{thm}
\begin{proof}
Let $L=\operatorname{colim}_k\Phi_kX$. Since $D$ is a finite $M$-set,
$\Hom_M(D,-)$ commutes with filtered colimits, so
$L^{+D}=\operatorname{colim}_k(\Phi_kX)^{+D}=\operatorname{colim}_{D_{k+1}=D}\Phi_{k+1}X=L$
with the unit of $L$ the identity of the colimit; thus $L$ is a $D$-sheaf,
likewise an $E$-sheaf, hence an $I'$-sheaf. The maps $X\to\Phi_kX$ are composites of maps of
$W_D\cup W_E$, hence lie in $W_D\vee W_E\subseteq W_{I'}$, and $a_{I'}=\Hom(P,-)$ commutes with filtered
colimits, so $a_{I'}L=a_{I'}X$ and $L=a_{I'}L=a_{I'}X$; no finiteness of $X$ is
used. Restricting to finite $X$, the functors
$\operatorname{colim}_k\Hom(T_k,-)$ and $\Hom(P,-)$ on finite $M$-sets are naturally
isomorphic, and by the Yoneda lemma for pro-objects $(T_k)_k\cong P$ in
$\mathrm{Pro}(\mathrm{Fin}M)$. Explicitly, the isomorphism consists of compatible maps
$b_k\colon P\to T_k$ and a map $a\colon T_{k_0}\to P$ with $ab_{k_0}=\mathrm{id}_P$
and, for each $j$ and all $k\gg j$, $b_j\,a\,\mu^k_{k_0}=\mu^k_j\colon T_k\to T_j$;
applying $\Hom(-,X)$ gives the retraction $\rho_k=b_k^*$ of
$\sigma_k=(a\mu^k_{k_0})^*\colon a_{I'}X\to\Phi_kX$ and the factorization of
$\Phi_{k_0}X\to\Phi_kX$ through $a_{I'}X$. Finally the image of
$T_{k+l}=D_{k+l}\otimes\cdots\otimes D_{k+1}\otimes T_k$ in $T_k$ is
$(D_{k+l}\cdots D_{k+1})\cdot T_k=I'\cdot T_k$ for $l\ge2m$, since the product of
$2m$ alternating ideals is $I'$; if $T_k=I'T_k\subseteq D_{k+1}T_k$ then
$\mu_{k+1}$ is surjective, the sizes $|T_k|$ are eventually non-increasing, hence
constant, and $\mu_k$ is eventually bijective; conversely eventual bijectivity
gives $T_k=\mathrm{im}(T_{k+2m}\to T_k)=I'T_k$. In that case
Lemma~\ref{lem:reduction} gives $W_D\vee W_E=W_{I'}$.
\end{proof}

\begin{cor}[finite categories]\label{cor:essconst}
Theorem~\ref{thm:essconst} holds for a finite category $S$ in place of $M$,
with $D,E$ idempotent two-sided ideals of $S$, $\otimes_M$ replaced by the
tensor product of $S$--$S$-bimodules (functors $S^{\mathrm{op}}\times S\to\mathbf{Set}$),
and $\Hom_M(T,X)$ replaced by the presheaf $c\mapsto\Hom_{\Presh(S)}(T(-,c),X)$.
\end{cor}
\begin{proof}
The ideals $D,E,I'$ and the tensors $T_k$, $P=I'\otimes_SI'$ are finite bimodules
(finitely many morphisms in each value), $X^{+D}(c)=\Hom(D(-,c),X)$ is the plus
construction for the covering sieve $D(-,c)$, and $a_{I'}X=\Hom(P,X)$ by
\cite[Thm.~6.10]{Topjoin}. Each $D(-,c)$ and $T_k(-,c)$ is a finite presheaf, so
$\Hom(D(-,c),-)$ and $\Hom(T_k(-,c),-)$ commute with filtered colimits, and the
first paragraph of the proof applies verbatim, objectwise in $c$ and for every
presheaf $X$: $L=\operatorname{colim}_k\Phi_kX$ is a $D$-sheaf and an $E$-sheaf, and
$L\cong a_{I'}X$. The functor $T\mapsto\Hom(T,-)$
from finite bimodules to functors $\mathrm{Fin}\Presh(S)\to\Presh(S)$ is fully
faithful: for each $c$ the functor $X\mapsto\Hom(T(-,c),X)$ on finite presheaves
is represented by the finite presheaf $T(-,c)$, so the Yoneda lemma on
$\mathrm{Fin}\Presh(S)$ recovers $T(-,c)$ together with the maps between such
functors, and naturality in $c$ recovers the bimodule structure. Hence the
Yoneda lemma for pro-objects applies in
$\mathrm{Pro}(\text{finite bimodules})$ and gives $(T_k)_k\cong P$ with the same
retract structure. Finally the product of $l\ge2m$ consecutive alternating
ideals is $I'$, so $\mathrm{im}(T_{k+l}\to T_k)=I'\cdot T_k$ as before, and
Lemma~\ref{lem:reduction} holds for presheaves on $S$ with the same proof.
\end{proof}

\subsection{The tower need not stabilize.}

\begin{prop}\label{prop:cyclic}
Let $\mathcal C$ be the category with two objects $a,b$, arrows $u\colon a\to b$,
$v\colon b\to a$, freely generated subject to the identification of all
composites of length $\ge3$ with the same endpoints; it has $10$ arrows. Let
$M$ be its contracted monoid ($12$ elements), $D$ the ideal of arrows factoring
through $b$ and $E$ that through $a$. Then $T_k\ne I'\cdot T_k$ for every even
$k$: the alternating plus-construction tower of $(D,E)$ never stabilizes.
\end{prop}
\begin{proof}
The identification of long words is a congruence because lengths add under
composition, so $\mathcal C$ is a category with, for each ordered pair of
objects, the words of length $\le2$ between them and one long arrow; $D$ and $E$
are idempotent (they contain identities), proper and incomparable, and
$I'$ is the ideal of long arrows and $0$. By Theorem~\ref{thm:essconst} it
suffices to exhibit, for every even $k$, a class of $T_k$ with no representative
whose first slot lies in $I'$. Take $k=2j$ and
$t=(1_a,uv,1_a,uv,\dots,1_a,uv)\in T_k$, whose slots alternate between $E$ and
$D$ as required. Every tuple of short arrows with product the word $w=(uv)^j$ of
the free category is a \emph{cut configuration} of $w$: cut points
$0=c_0\le c_1\le\dots\le c_k=2j$, the $i$-th slot being the subword between
$c_{i-1}$ and $c_i$, and an empty slot being the identity of the object at its
cut point, which must be $a$ (even cut point) for an $E$-slot ($i$ odd) and $b$
(odd cut point) for a $D$-slot ($i$ even). A tensor move between two tuples of
short arrows moves one cut point and preserves $w$. Put $d_i:=i-c_i$, so
$d_0=d_k=0$, a slot of length $\ell$ changes $d$ by $1-\ell$, and an empty slot
raises $d$ by one and is allowed only when $d_i$ is odd, i.e.\ only from an even
value of $d$. From an odd value $d$ cannot rise, since the next slot must be
non-empty; hence $d\le1$ throughout, and once $d\le-1$ it can never return to
$0$ (from an even value it rises by one to an odd value and is then stuck).
A slot of length $\ge3$ would lower $d$ from at most $1$ to at most $-1$, so no
cut configuration of $w$ has a slot of length $\ge3$. Consequently no tensor
move applied to a cut configuration produces a long slot (the target would be a
cut configuration with a slot of length $\ge3$) or a zero slot (adjacent slots
of a configuration are composable), and the class of $t$ consists of cut
configurations only; its first slot is never in $I'$.
\end{proof}

\begin{rem}
Enumerating the tensor relations gives $|T_1|=10$ and $|T_k|=9$ for $k\ge2$,
with $|I'T_k|=5$: four classes of short tuples survive at every level. For the
analogous monoid with three objects in a cycle and words of length $\ge4$
identified ($|M|=23$) the sizes oscillate, $19,19,16,19,16,\dots$.
\end{rem}

Whether $W_D\vee W_E=W_{J_D\vee J_E}$ holds on the monoids of
Proposition~\ref{prop:cyclic} is not decided by Lemma~\ref{lem:reduction}, which
is only sufficient; by Theorem~\ref{thm:essconst} no pro-limit of the tower can
obstruct it, and the dual tower $\lim_kX\otimes_M(D\otimes D)\otimes(E\otimes E)\otimes\cdots$,
which inverts $W_D$ and $W_E$, is $X\otimes P$ and does not obstruct it either.

\begin{que}\label{que:finite}
Is $\Top(S)$ a sublattice of $\Same(\Presh(S))$ for every finite category $S$?
Equivalently, is the unit $X\to a_{J\vee J'}X$ in $W_J\vee W_{J'}$ for all $X$?
The first undecided cases are the monoids of Proposition~\ref{prop:cyclic}.
\end{que}

\begin{que}
Classify the natural samenesses (Proposition~\ref{prop:natural}); is the
interval $[\Iso,\Wcyc]$ of natural samenesses finite?
\end{que}

\begin{que}
Does Theorem~\ref{thm:I} hold for cubical objects of an $\infty$-topos with the
cubical coskeleta? For cubical sets it does \cite{Postnikov}; the missing
ingredient is a cubical form of Proposition~\ref{prop:const}.
\end{que}

\section*{Funding}
No funding was received for this work.

\section*{Declarations}
\noindent\textbf{CRediT authorship contribution statement.}\quad Ryota Shibusawa:
Conceptualization, Methodology, Formal analysis, Investigation, Software,
Validation, Visualization, Writing -- original draft, Writing -- review \& editing.

\noindent\textbf{Competing interests.}\quad The author declares that he has no
competing interests.

\begin{thebibliography}{99}
\bibitem{ABFJ} M.~Anel, G.~Biedermann, E.~Finster, A.~Joyal, \emph{Left-exact
localizations of $\infty$-topoi I: Higher sheaves}, Adv. Math. \textbf{400}
(2022), 108268.
\bibitem{ABFJ2} M.~Anel, G.~Biedermann, E.~Finster, A.~Joyal, \emph{Left-exact
localizations of $\infty$-topoi II: Grothendieck topologies}, J. Pure Appl.
Algebra \textbf{228} (2024), 107597.
\bibitem{Duskin} J.~Duskin, \emph{Simplicial methods and the interpretation
of ``triple'' cohomology}, Mem. Amer. Math. Soc. \textbf{163} (1975).
\bibitem{GJ} P.~G.~Goerss, J.~F.~Jardine, \emph{Simplicial Homotopy Theory},
Progress in Mathematics \textbf{174}, Birkh\"auser, 1999.
\bibitem{HN} P.~Hell, J.~Ne\v{s}et\v{r}il, \emph{Graphs and Homomorphisms},
Oxford Lecture Series in Mathematics and its Applications \textbf{28}, Oxford
University Press, 2004.
\bibitem{Elephant} P.~T.~Johnstone, \emph{Sketches of an Elephant: A Topos
Theory Compendium}, Oxford University Press, 2002.
\bibitem{Lueck} W.~L\"uck, \emph{Transformation Groups and Algebraic $K$-Theory},
Lecture Notes in Mathematics \textbf{1408}, Springer, 1989.
\bibitem{KS} M.~Kashiwara, P.~Schapira, \emph{Categories and Sheaves},
Grundlehren der mathematischen Wissenschaften \textbf{332}, Springer, 2006.
\bibitem{KL89} G.~M.~Kelly, F.~W.~Lawvere, \emph{On the complete lattice of
essential localizations}, Bull. Soc. Math. Belg. S\'er. A \textbf{41} (1989),
289--319.
\bibitem{KRRZ} C.~Kennett, E.~Riehl, M.~Roy, M.~Zaks, \emph{Levels in the
toposes of simplicial sets and cubical sets}, J. Pure Appl. Algebra
\textbf{215} (2011), 949--961.
\bibitem{HTT} J.~Lurie, \emph{Higher Topos Theory}, Annals of Mathematics
Studies \textbf{170}, Princeton University Press, 2009.
\bibitem{Marques25} J.~Marqu\`es, \emph{A criterion for categories on which
every Grothendieck topology is rigid}, Appl. Categ. Structures \textbf{33}
(2025), art.\ 37.
\bibitem{MLM} S.~Mac Lane, I.~Moerdijk, \emph{Sheaves in Geometry and Logic},
Springer, 1992.
\bibitem{Same} R.~Shibusawa, \emph{The sameness lattice of a category:
weak-equivalence structures, Galois functoriality, and comparison intervals},
preprint (2026), archived at DOI 10.5281/zenodo.22823772.
\bibitem{Postnikov} R.~Shibusawa, \emph{Postnikov truncations as joins of sheaf
and homotopy equivalences}, preprint (2026), archived at DOI 10.5281/zenodo.22956181.
\bibitem{Topjoin} R.~Shibusawa, \emph{Meets and joins of Grothendieck
topologies in the sameness lattice}, preprint (2026), archived at DOI
10.5281/zenodo.22978446.
\bibitem{EI} R.~Shibusawa, \emph{The EI-reflection of a category, its preorder
reflection, and the homomorphism order of graphs}, preprint (2026), archived at
DOI 10.5281/zenodo.23029859.
\bibitem{Infjoin} R.~Shibusawa, \emph{Postnikov truncation of simplicial objects
in an $\infty$-topos as a join of localizations}, preprint (2026), archived at
DOI 10.5281/zenodo.23049588.
\end{thebibliography}

\end{document}
